\chapter{Control theory and the LIC101 loop}
\label{ch:control}

This chapter explains feedback control from first principles and then the
specific controller in \code{pyplant/plc/pid.py}: a velocity-form PID with
structural anti-windup, bumpless transfer, derivative on filtered measurement
and IMC tuning.

\section{Feedback control}

A feedback loop measures the controlled variable (the \emph{process value},
PV), compares it with the desired value (the \emph{setpoint}, SP), and moves
a manipulated variable (the \emph{controller output}, OP) to remove the
difference $e = \text{SP} - \text{PV}$. Figure~\ref{fig:loop} shows LIC101 in
this form.

\begin{figure}[htbp]
\centering
\scalebox{0.8}{%
\begin{tikzpicture}[node distance=11mm]
  \node (sp) {SP};
  \node[circle,draw,right=8mm of sp,inner sep=2pt] (sum) {$\Sigma$};
  \node[hot,right=of sum] (c) {PID $C(s)$\\LIC101};
  \node[box,right=of c] (zoh) {ZOH\\$T_s=0.1$\,s};
  \node[box,right=of zoh] (p) {valve + tank\\$G(s)=\dfrac{K}{\tau s+1}$};
  \node[circle,draw,right=10mm of p,inner sep=2pt] (dsum) {$\Sigma$};
  \node[right=8mm of dsum] (pv) {level};
  \node[above=8mm of dsum] (d) {disturbance (leak)};
  \node[box,below=10mm of zoh] (m) {LT101\\+ noise};
  \draw[flow] (sp) -- (sum);
  \draw[flow] (sum) -- node[lbl,above]{$e$} (c);
  \draw[flow] (c) -- node[lbl,above]{OP} (zoh);
  \draw[flow] (zoh) -- (p);
  \draw[flow] (p) -- (dsum);
  \draw[flow] (d) -- (dsum);
  \draw[flow] (dsum) -- (pv);
  \draw[flow] ($(dsum.east)+(4mm,0)$) |- (m);
  \draw[flow] (m) -| node[lbl,pos=0.9,left]{$-$PV} (sum);
\end{tikzpicture}}
\caption{LIC101 as a feedback loop. The PLC samples the level every 100\,ms and
holds its output between scans.}
\label{fig:loop}
\end{figure}

\section{The PID law}

The ideal continuous PID controller in standard (ISA) form is
\begin{equation}
  u(t) = K_c\left( e(t) + \frac{1}{T_i}\int_0^t e(\tau)\,d\tau
         + T_d\frac{de}{dt}\right),
\end{equation}
with controller gain $K_c$, integral time $T_i$ and derivative time $T_d$. Each
term has a role:
\begin{description}
  \item[Proportional] reacts to the present error. Alone it leaves a steady
        offset whenever a non-zero output is needed to hold the setpoint.
  \item[Integral] accumulates past error and keeps moving the output until
        the error is exactly zero. It removes offset, and it is the reason a
        leak is masked (Section~\ref{sec:masking}).
  \item[Derivative] anticipates by reacting to the rate of change. On a noisy
        level it does more harm than good, so LIC101 runs with $T_d=0$.
\end{description}

\subsection{Positional and velocity forms}

Discretising at sample time $T_s$ with backward differences gives the
\emph{positional} form
\begin{equation}
  u_k = K_c\Big(e_k + \frac{T_s}{T_i}\sum_{j\le k} e_j
        + \frac{T_d}{T_s}(e_k - e_{k-1})\Big).
\end{equation}
Subtracting $u_{k-1}$ from $u_k$ gives the \emph{velocity} (incremental) form
used here:
\begin{equation}
  \Delta u_k = K_c(e_k - e_{k-1}) + K_c\frac{T_s}{T_i}e_k
   - K_c\frac{T_d}{T_s}\big(y^f_k - 2y^f_{k-1} + y^f_{k-2}\big),
  \qquad u_k = \operatorname{clip}\big(u_{k-1} + \Delta u_k,\,0,\,100\big),
  \label{eq:velocity}
\end{equation}
where $y^f$ is the filtered measurement (below). For a reverse-acting loop the
derivative term has a minus sign because it acts on the measurement.

\subsection{Why the velocity form}

In the velocity form the output itself is the integrator state. Two
properties that need extra logic in the positional form become automatic:

\begin{theory}[Anti-windup is structural]
When the actuator saturates, a positional PID keeps integrating the error, so
the integral grows far beyond the limit (\emph{windup}) and the output stays
pinned long after the error reverses. In equation~\eqref{eq:velocity} the
increment that would push past a limit is simply clipped away. No hidden
accumulator grows, and the output leaves the limit on the first scan the
error changes sign.
\end{theory}

\begin{theory}[Bumpless transfer is structural]
In manual mode the output is whatever the operator writes, and the controller
keeps its error history fresh against the current PV (\code{track()}). On
return to auto it continues incrementing from the manual value, so the valve
does not jump.
\end{theory}

\subsection{Derivative on filtered measurement}

Taking the derivative of the error differentiates the setpoint too, so a
setpoint step would produce an impulse (\emph{derivative kick}). The code
differentiates the measurement instead. The measurement is first passed
through a first-order low-pass filter with time constant $T_f = 5$\,s,
discretised as
\begin{equation}
  y^f_k = y^f_{k-1} + \alpha\,(y_k - y^f_{k-1}),
  \qquad \alpha = \frac{T_s}{T_f + T_s}.
\end{equation}

\subsection{Initialisation from valve feedback}

A velocity PID increments from wherever its output already is. Starting at
0\% would slam shut the valve of a plant running at 66\%, dead-head the pump
and trip the interlock on the first scan. On the first scan the program
therefore seeds the output from the measured stem position \tg{ZT101}.

\section{IMC tuning}
\label{sec:imc}

Internal Model Control (IMC) tuning chooses the controller so that the closed
loop behaves like a first-order lag with a time constant $\lambda$ the
engineer specifies. For a first-order process
\begin{equation}
  G(s) = \frac{K}{\tau s + 1},
\end{equation}
the IMC controller is $Q(s) = G^{-1}(s)\,\frac{1}{\lambda s+1}$ and the
equivalent feedback controller is
\begin{equation}
  C(s) = \frac{Q}{1 - GQ} = \frac{\tau s + 1}{K\lambda s}
       = \frac{\tau}{K\lambda}\left(1 + \frac{1}{\tau s}\right),
\end{equation}
which is a PI controller with
\begin{equation}
  \boxed{K_c = \frac{\tau}{K\lambda}, \qquad T_i = \tau .}
\end{equation}
The integral time cancels the process pole, and the closed-loop response to a
setpoint step is
\begin{equation}
  \frac{\text{PV}(s)}{\text{SP}(s)} = \frac{GC}{1+GC} = \frac{1}{\lambda s+1}.
\end{equation}

For LIC101, $K = 15.4$\,mm/\%, $\tau = 72$\,s and $\lambda = 13$\,s:
\[
  K_c = \frac{72}{15.4\times13} \approx 0.36 \;\to\; 0.35\ \%/\text{mm},
  \qquad T_i = 72\ \text{s}, \qquad T_d = 0.
\]
The ratio $\tau/\lambda\approx5.5$ gives a closed loop five times faster than
the open loop while leaving margin for the neglected valve lag, the sample
time and the nonlinearity.

\begin{figure}[htbp]
\centering
\begin{tikzpicture}
\begin{axis}[width=0.6\textwidth,height=4.2cm,xlabel={time after step (s)},
  ylabel={level (mm)},grid=major,grid style={linegrey},xmin=0,xmax=120,
  ymin=290,ymax=360,legend pos=south east]
\addplot[slate,dashed,thick] coordinates {(0,300)(0,350)(120,350)};
\addplot[accent,very thick,domain=0:120,samples=120]{300+50*(1-exp(-x/13))};
\addplot[badred,thick,domain=0:120,samples=120]{300+50*(1-exp(-x/72))};
\addplot[only marks,mark=|,mark size=4pt] coordinates {(52,349)};
\legend{setpoint, closed loop ($\lambda=13$\,s), open loop ($\tau=72$\,s), 2\% settled}
\end{axis}
\end{tikzpicture}
\caption{Ideal response to a 50\,mm setpoint step. The measured response of
the real loop matches: no overshoot, 2\% settling in about 40--50\,s, zero
offset.}
\label{fig:step}
\end{figure}

\section{Sampling and the zero-order hold}

The PLC executes every $T_s = 100$\,ms and holds its output until the next
scan. A zero-order hold behaves approximately like a dead time of $T_s/2 =
50$\,ms, negligible against $\lambda = 13$\,s, but it is modelled physically
rather than ignored: the simulator integrates at 50\,ms with the input held
between scans. A PID that looks better in simulation than on a PLC usually
lacks exactly this.

\section{Why integral action hides a leak}
\label{sec:masking}

A leak is a load disturbance $d$ on tank~1. With a PI controller the
steady-state error to a step disturbance is, by the final value theorem,
\begin{equation}
  e_\infty = \lim_{s\to0} s\cdot\frac{G_d(s)}{1+G(s)C(s)}\cdot\frac{d}{s} = 0,
\end{equation}
because $C(s)$ has a pole at $s=0$. The controller drives the level back to
setpoint exactly, and the only lasting evidence is that the output has moved:

\begin{center}
\begin{tabular}{@{}lrr@{}}
\toprule
& \tg{LT101} & \tg{LIC101_OP} \\
\midrule
before the leak & 300.00\,mm & 66.18\% \\
after the leak & 300.00\,mm & 75.80\% \\
\bottomrule
\end{tabular}
\end{center}

A diagnosis model fed only the process value cannot see this fault. This is
why the detectors are given the controller output, the valve position and
derived physics channels (Chapter~\ref{ch:ml}).

\section{Controller parameters}

\begin{center}
\begin{tabular}{@{}llll@{}}
\toprule
\code{PidTuning} field & Default & Unit & Meaning \\
\midrule
\code{kc} & 0.35 & \%/mm & controller gain \\
\code{ti} & 72.0 & s & integral time (0 disables integral action) \\
\code{td} & 0.0 & s & derivative time (0 disables derivative) \\
\code{tf} & 5.0 & s & derivative filter time constant \\
\code{out_lo}, \code{out_hi} & 0, 100 & \% & output limits \\
\code{direct_acting} & False & & False: output rises when PV falls \\
\bottomrule
\end{tabular}
\end{center}

Chapter~\ref{ch:params} explains how to change them.
